\documentclass{article}
\usepackage{graphicx} % Required for inserting images
\usepackage{booktabs}

\title{resultados P3\_electronica}
\author{Jesús Brandon Villanueva Rivas}
\date{September 2026}

\begin{document}

\maketitle


\section{Resultados}

\begin{table}[htbp]
\caption{Caracterización del diodo 1N4007.}
\label{tab:dato_1N4007}
\begin{tabular}{cccc}
\toprule
$V_S$ [V] & $V_D$ [V] & $V_R$ [V] & $I_D$ [mA] \\
\midrule
0.0 & 0.0000 & 0.000 & 0.00 \\
0.1 & 0.1465 & 0.000 & 0.00 \\
0.2 & 0.2149 & 0.000 & 0.00 \\
0.3 & 0.3102 & 0.000 & 0.00 \\
0.4 & 0.4170 & 0.000 & 0.01 \\
0.5 & 0.4900 & 0.010 & 0.13 \\
0.6 & 0.5450 & 0.055 & 0.53 \\
0.7 & 0.5700 & 0.130 & 0.73 \\
0.8 & 0.5880 & 0.212 & 1.02 \\
0.9 & 0.6130 & 0.287 & 1.33 \\
1.0 & 0.6210 & 0.379 & 1.87 \\
1.1 & 0.6280 & 0.472 & 2.12 \\
1.2 & 0.6390 & 0.561 & 2.62 \\
1.3 & 0.6450 & 0.655 & 3.12 \\
1.4 & 0.6510 & 0.749 & 3.56 \\
1.5 & 0.6590 & 0.841 & 3.92 \\
1.6 & 0.6620 & 0.938 & 4.16 \\
1.7 & 0.6670 & 1.033 & 4.78 \\
1.8 & 0.6700 & 1.130 & 5.13 \\
1.9 & 0.6740 & 1.226 & 5.51 \\
2.0 & 0.6770 & 1.323 & 5.98 \\
2.1 & 0.6800 & 1.420 & 6.29 \\
2.2 & 0.6830 & 1.517 & 6.77 \\
2.3 & 0.6860 & 1.614 & 7.18 \\
2.4 & 0.6890 & 1.711 & 7.68 \\
2.5 & 0.6910 & 1.809 & 8.06 \\
2.6 & 0.6950 & 1.905 & 8.41 \\
2.7 & 0.6970 & 2.003 & 9.14 \\
2.8 & 0.6980 & 2.102 & 9.65 \\
2.9 & 0.7000 & 2.200 & 9.94 \\
3.0 & 0.7020 & 2.298 & 10.34 \\
3.1 & 0.7020 & 2.398 & 10.64 \\
3.2 & 0.7030 & 2.497 & 11.04 \\
3.3 & 0.7050 & 2.595 & 11.43 \\
3.4 & 0.7070 & 2.693 & 12.04 \\
3.5 & 0.7080 & 2.792 & 12.31 \\
3.6 & 0.7100 & 2.890 & 12.73 \\
3.7 & 0.7120 & 2.988 & 13.39 \\
3.8 & 0.7130 & 3.087 & 13.79 \\
3.9 & 0.7140 & 3.186 & 14.13 \\
4.0 & 0.7160 & 3.284 & 14.50 \\
4.1 & 0.7170 & 3.383 & 15.09 \\
4.2 & 0.7190 & 3.481 & 15.51 \\
4.3 & 0.7200 & 3.580 & 15.92 \\
4.4 & 0.7210 & 3.679 & 16.33 \\
4.5 & 0.7220 & 3.778 & 16.71 \\
4.6 & 0.7240 & 3.876 & 17.44 \\
4.7 & 0.7240 & 3.976 & 17.80 \\
4.8 & 0.7250 & 4.075 & 18.22 \\
4.9 & 0.7270 & 4.173 & 18.84 \\
5.0 & 0.7280 & 4.272 & 19.30 \\
\bottomrule
\end{tabular}
\end{table}



Dados los datos de la tabla \ref{tab:dato_1N4007} y los datos simulados en Kicad se construyo la gráfica \ref{} y \ref{}.
La gráfica \ref{fig:VD_VS} muestra cómo cambia el voltaje en las terminales del diodo ($V_D$) respecto al voltaje de entrada ($V_S$), donde se aprecian dos regiones principales:

\begin{itemize}
  \item Cuando $V_S < 0.6\text{ V}$, prácticamente toda la caída de tensión se concentra en el diodo ($V_D \approx V_S$) y el voltaje en la resistencia es nulo.
  \item Cuando $V_S > 0.6\text{ V}$, el voltaje en el diodo se estabiliza (aumenta muy suavemente) y el resto de la caída de tensión se transfiere a la resistencia.
\end{itemize}



De igual manera, se pueden observar tres regiones principales en la gráfica \ref{fig:ID_VD}:
\begin{itemize}
  \item Cuando $V_D < 0.4900~\text{V}$, la corriente es prácticamente nula (por debajo de la resolución del multímetro).
  \item En el rango $0.4900~\text{V} \le V_D \le 0.5880~\text{V}$ comienza el encendido del diodo (inicio de la conducción). Además, a partir de los datos de la tabla se aprecia que en este intervalo ocurre un incremento de un orden de magnitud en la corriente (de $0.13~\text{mA}$ a $1.02~\text{mA}$).
  \item Cuando $V_D > 0.5880~\text{V}$, la corriente en el diodo crece de manera exponencial, entrando completamente en la zona de conducción.
\end{itemize}

Notemos que la zona de conducción inicia antes de los $0.7~\text{V}$. Este valor es teórico y proviene de un modelo simplificado (aproximación de caída constante). En esta práctica, el diodo comenzó a conducir de forma medible desde los $0.49~\text{V}$.

Asimismo, la comparación entre los datos experimentales y simulados deja ver un desplazamiento hacia la izquierda de la curva experimental respecto a la simulada. 
La simulación y las hojas de datos consideran un entorno a temperatura de referencia ($25~^{\circ}\text{C}$); no obstante, en la práctica la potencia disipada por el propio diodo genera calor por efecto Joule. Este incremento térmico reduce la barrera de potencial del silicio, adelantando el codo de conducción. 
Otro factor no contemplado en el modelo genérico de la simulación son las tolerancias de fabricación del componente real, donde la concentración de dopado puede variar e introducir desviaciones en sus parámetros físicos (como la corriente de saturación inversa $I_s$ y el factor de idealidad $\eta$). 
Finalmente, la resolución del multímetro y el efecto de carga introducido por el multímetro al medir corriente en serie afectan la lectura.

\end{document}
